\section*{Chapitre \ref{ch:g7:earthquakes} --- Les séismes}

\begin{solution}{exo:g7:earthquakes:1}
(a) \omterm{def:g7:earthquakes:focus}{foyer}~; (b) \omterm{def:g7:earthquakes:intensity}{intensité}~; (c) \omterm{def:g7:earthquakes:fault}{faille}~; (d) \omterm{def:g7:earthquakes:focus}{épicentre}~; (e) \omterm{def:g7:earthquakes:magnitude}{magnitude}~;
(f) \omterm{def:g7:earthquakes:seismometer}{sismomètre}.
\end{solution}

\begin{solution}{exo:g7:earthquakes:2}
(a) Au village lui-même~: l'\omterm{def:g7:earthquakes:focus}{épicentre} est le point de la surface situé
juste au-dessus du \omterm{def:g7:earthquakes:focus}{foyer}. (b) \qty{30}{km}, à la verticale. (c) En général
plus faible~: la secousse s'affaiblit avec la distance. (Cela dépend aussi
du sol~: sur de la \omterm{def:g3:stones-pebbles-sand:sand}{boue} molle, le second village pourrait tout de même
trembler fort.)
\end{solution}

\begin{solution}{exo:g7:earthquakes:3}
(a) \omterm{def:g7:earthquakes:magnitude}{Magnitude} (un nombre avec une décimale). (b) \omterm{def:g7:earthquakes:intensity}{Intensité} (ce qui a été
vu en un lieu). (c) \omterm{def:g7:earthquakes:intensity}{Intensité} (chiffres romains, différents en deux
lieux). (d) \omterm{def:g7:earthquakes:magnitude}{Magnitude}~: l'énergie libérée, un pas de \omterm{def:g7:earthquakes:magnitude}{magnitude} de plus.
\end{solution}

\begin{solution}{exo:g7:earthquakes:4}
(a) Faux~: les scientifiques ne savent pas prévoir les \omterm{def:g4:volcanoes-earthquakes-first-look:earthquake}{séismes}~; ils
savent seulement quelles \omterm{def:g7:earthquakes:fault}{failles} rompront de nouveau un jour. (b) Faux~:
en plein \omterm{def:g2:water-on-earth:ocean}{océan}, un \omterm{def:g7:earthquakes:tsunami}{tsunami} dépasse rarement un mètre de haut~; il grandit
près de la \omterm{def:g1:landscapes-around-us:water}{côte}. (c) Vrai. (d) Vrai~: 92~\% des \omterm{def:g4:volcanoes-earthquakes-first-look:earthquake}{séismes} de \omterm{def:g7:earthquakes:magnitude}{magnitude}~6 ou
plus se produisent à moins de \qty{300}{km} d'une \omterm{def:g7:plate-tectonics:plate-boundary}{limite de plaques}.
\end{solution}

\begin{solution}{exo:g7:earthquakes:5}
(a) Environ 127 par an de \omterm{def:g7:earthquakes:magnitude}{magnitude} 6 à 6.9~; environ 13 de \omterm{def:g7:earthquakes:magnitude}{magnitude} 7 à
7.9. (b) $126.7 \div 13.12 \approx 9.7$~: environ 10~fois moins. (c)
$13.12 + 0.76 + 0.04 = 13.92$~: environ \textbf{14} par an. (d)
$0.04 \times 50 = 2$~: deux en 50~ans (Sumatra en 2004 et le Japon en
2011).
\end{solution}

\begin{solution}{exo:g7:earthquakes:6}
(a) Un pas~: 32~fois. (b) Deux pas~: $32 \times 32 = 1024$, environ
1000~fois. (c) Chaque pas rend les ondes 10~fois plus grandes~: $10 \times 10 =
100$~fois plus grandes.
\end{solution}

\begin{solution}{exo:g7:earthquakes:7}
(a) $8 \times 5 = \qty{40}{km}$. (b) $8 \times 25 = \qty{200}{km}$. (c)
$120 \div 8 = \qty{15}{s}$. (d) Une seule station donne seulement une
distance, pas une direction~: tous les points situés à cette distance
forment un cercle autour de la station, et l'\omterm{def:g7:earthquakes:focus}{épicentre} pourrait se trouver
n'importe où sur ce cercle.
\end{solution}

\begin{solution}{exo:g7:earthquakes:8}
(a) V. (b) II. (c) VIII. (d) L'\omterm{def:g7:earthquakes:intensity}{intensité} change d'un lieu à l'autre~: le
témoignage (c) vient d'un lieu proche de l'\omterm{def:g7:earthquakes:focus}{épicentre}, ou bâti sur un sol
mou~; le témoignage (a), d'un lieu plus éloigné, ou bâti sur la \omterm{def:g3:stones-pebbles-sand:rock}{roche}
dure. Le \omterm{def:g4:volcanoes-earthquakes-first-look:earthquake}{séisme} n'a toujours qu'une seule \omterm{def:g7:earthquakes:magnitude}{magnitude}.
\end{solution}

\begin{solution}{exo:g7:earthquakes:9}
(a) $7031 - 5561 = 1470$. (b) $1470 \div 7031 \approx 0.209$~: environ
\textbf{21~\%}. (c) La \omterm{def:g7:plate-tectonics:plate}{plaque} qui s'enfonce descend profondément dans le
\omterm{def:g6:inside-the-earth:layers}{manteau}. Elle est encore assez rigide pour se rompre~: des \omterm{def:g4:volcanoes-earthquakes-first-look:earthquake}{séismes} peuvent
donc commencer à l'intérieur, à des centaines de kilomètres de profondeur.
\end{solution}

\begin{solution}{exo:g7:earthquakes:10}
(a) $2500 \div 800 = 3.125$~heures~: 3~heures et $0.125 \times 60 = 7.5$
minutes, soit environ 3~heures 8~minutes. (b) 3~h~8~min $-$ 2~h~17~min $=$
51~minutes plus tard. (c) La vague a mis des heures à arriver, alors qu'un
grand \omterm{def:g4:volcanoes-earthquakes-first-look:earthquake}{séisme} sous la \omterm{def:g2:water-on-earth:ocean}{mer} est localisé en quelques minutes~: une alerte
aurait pu envoyer les gens sur les hauteurs bien avant l'arrivée de la
vague.
\end{solution}

\begin{solution}{exo:g7:earthquakes:11}
(a) $400 \div 800 = 0.5$~heure~: 30~minutes. (b) $20 \div 40 = 0.5$~heure~:
30~minutes. (c) 60~minutes. (d) Les \qty{20}{km} de la fin prennent autant
de temps que les \qty{400}{km} du début~: la vague ralentit énormément en
eau peu profonde. (e) L'avant de la vague ralentit tandis que l'eau qui le
suit arrive encore vite~: l'eau s'entasse, et la vague grandit.
\end{solution}

\begin{solution}{exo:g7:earthquakes:12}
(a) Trois pas~: $32 \times 32 \times 32 = 32\,768$, environ
\textbf{33\,000}~fois. (b) $32\,768 \div 127 \approx 258$~: environ
260~années de tous les \omterm{def:g4:volcanoes-earthquakes-first-look:earthquake}{séismes} de \omterm{def:g7:earthquakes:magnitude}{magnitude}~6 du monde. (c) Presque toute
l'énergie est libérée par les quelques très grands \omterm{def:g4:volcanoes-earthquakes-first-look:earthquake}{séismes}~: un seul
\omterm{def:g4:volcanoes-earthquakes-first-look:earthquake}{séisme} de \omterm{def:g7:earthquakes:magnitude}{magnitude}~9 pèse plus lourd que des siècles de \omterm{def:g4:volcanoes-earthquakes-first-look:earthquake}{séismes} de
\omterm{def:g7:earthquakes:magnitude}{magnitude}~6.
\end{solution}

\begin{solution}{pb:g7:earthquakes:1}
\textbf{1.} A~: $45 - 25 = 20$~s. B~: $29 - 16 = 13$~s. C~: $34 - 19 =
15$~s.

\textbf{2.} La station B~: son décalage est le plus court, et plus le
décalage est court, plus le \omterm{def:g4:volcanoes-earthquakes-first-look:earthquake}{séisme} est proche. (Ses ondes rapides sont
aussi arrivées les premières.)

\textbf{3.} A~: $8 \times 20 = \qty{160}{km}$. B~: $8 \times 13 =
\qty{104}{km}$. C~: $8 \times 15 = \qty{120}{km}$.

\textbf{4.} Tous les points situés à \qty{160}{km} de A~: un cercle de
rayon \qty{160}{km} centré sur A. L'\omterm{def:g7:earthquakes:focus}{épicentre} se trouve quelque part
dessus.

\textbf{5.} A~: $160 \div 20 = 8$~carreaux. B~: $104 \div 20 = 5.2$~carreaux.
C~: $120 \div 20 = 6$~carreaux.

\textbf{6.} E est à $3 + 6 = 9$~carreaux à l'est et 14.5~carreaux au nord
de O, en \omterm{def:g2:water-on-earth:ocean}{mer}, entre les trois stations~; les cercles de A et de B passent
par ce point.

\textbf{7.} Sous la \omterm{def:g2:water-on-earth:ocean}{mer}. Un mouvement brusque du fond de la \omterm{def:g2:water-on-earth:ocean}{mer} peut
déclencher un \omterm{def:g7:earthquakes:tsunami}{tsunami}~: il faut alerter les \omterm{def:g1:landscapes-around-us:water}{côtes} et l'île.

\textbf{8.} Un pas de \omterm{def:g7:earthquakes:magnitude}{magnitude}~: environ 32~fois plus d'énergie.

\textbf{9.} $12 \times 20 = \qty{240}{km}$.

\textbf{10.} $240 \div 800 = 0.3$~heure, et $0.3 \times 60 = 18$~minutes.

\textbf{11.} 5~minutes après le \omterm{def:g4:volcanoes-earthquakes-first-look:earthquake}{séisme}.

\textbf{12.} $18 - 5 = 13$~: les habitants de Port Hale ont environ
\textbf{13~minutes} pour gagner les hauteurs. (La vague ralentit en eau
peu profonde près de l'île, ce qui laisse un peu plus de temps, mais
personne ne doit compter dessus~: ils partent tout de suite.)
\end{solution}
